\chapter{Structures, Type Classes, and Coercions}

% ADDED: Chapter orientation paragraph
In this chapter, we explore how mathematical objects and their
properties are represented in Lean and Mathlib.
We begin with polymorphism and implicit arguments, then examine how
Lean models rational numbers using structures
(Section~\ref{sec:exploring_mathlib}) and coercions
(Section~\ref{sec:coercions}).
Along the way, we introduce type classes as a mechanism for
overloading operations across different types.

It is interesting to note that a relation, in type theory, can be
expressed as a function:
\lstinline[language=lean]|R : α → α → Prop|.
Similarly, when defining a predicate
(\lstinline[language=lean]|P : α → Prop|) we must first declare
\lstinline[language=lean]|α : Type| to be some arbitrary type.
This is what is called \textbf{polymorphism}, more specifically
\textbf{parametric polymorphism}.
A canonical example is the identity function, written as
\lstinline[language=lean]|α → α|, where
\lstinline[language=lean]|α| is a type variable.
It has the same type for both its domain and codomain, meaning it can
be applied to booleans (returning a boolean), numbers (returning a
number), functions (returning a function), and so on.

In the same spirit, we can define a transitivity property of a
relation, for a given type \lstinline[language=lean]|α|, as follows:
\begin{lstlisting}[language=lean]
def Transitive (α : Type) (R : α → α → Prop) : Prop :=
  ∀ x y z, R x y → R y z → R x z
\end{lstlisting}
To use \lstinline[language=lean]|Transitive|, we must provide both
the type \lstinline[language=lean]|α| and the relation itself.
For example, here is a proof of transitivity for the
less-than-or-equal relation on $\mathbb{N}$
(in Lean, \lstinline[language=lean]|Nat| or
\lstinline[language=lean]|ℕ|):
\begin{lstlisting}[language=lean]
theorem le_trans_nat : Transitive Nat (· ≤ · : Nat → Nat → Prop) :=
  fun x y z h1 h2 => Nat.le_trans h1 h2
\end{lstlisting}
Notice that we used a function to discharge the universal quantifiers
required by transitivity.
Looking at this code, we immediately notice that explicitly passing
the type argument \lstinline[language=lean]|Nat| is somewhat
repetitive.
Lean allows us to omit it by letting the type inference mechanism
fill it in automatically.
This is achieved by using \textbf{implicit arguments} with curly
brackets:
\begin{lstlisting}[language=lean]
def Transitive {α : Type} (R : α → α → Prop) : Prop :=
  ∀ x y z, R x y → R y z → R x z
theorem le_trans_nat' : Transitive (· ≤ · : Nat → Nat → Prop) :=
  fun x y z h1 h2 => Nat.le_trans h1 h2
\end{lstlisting}

Let us now revisit the transitivity proof, but this time for the
less-than-or-equal relation on the rational numbers
(\lstinline[language=lean]|Rat| or \lstinline[language=lean]|ℚ|):
\begin{lstlisting}[language=lean]
import Mathlib

theorem rat_le_trans : Transitive (· ≤ · : Rat → Rat → Prop) :=
  fun _ _ _ h1 h2 => Rat.le_trans h1 h2
\end{lstlisting}
Here, \lstinline[language=lean]|Rat.le_trans| is the transitivity
lemma for \lstinline[language=lean]|≤| on rational numbers, provided
by Mathlib.
We import Mathlib to access \lstinline[language=lean]|Rat| and
\lstinline[language=lean]|le_trans|.
Mathlib is the community-driven mathematical library for Lean,
containing a large body of formalized mathematics under ongoing
development.
It is the de facto standard library for mathematical proofs in
Lean~(\cite{mathlib2020}).
We will explore it throughout this chapter.
The underscores indicate unnamed variables that we do not reference
later.
If we had named them, say \lstinline|x y z|, then
\lstinline[language=lean]|h1| would be a proof of
\lstinline[language=lean]|x ≤ y|,
\lstinline[language=lean]|h2| would be a proof of
\lstinline[language=lean]|y ≤ z|, and
\lstinline[language=lean]|Rat.le_trans h1 h2| produces a proof of
\lstinline[language=lean]|x ≤ z|.
The \lstinline[language=lean]|Transitive| definition is imported from
Mathlib and similarly defined as before.

\begin{example}
  The code can be made more readable using \textbf{tactic mode}.
  In this mode, you use tactics---commands provided by Lean or
  defined by users---to carry out proof steps succinctly, avoid
  repetition, and automate common patterns.
  This often yields shorter, clearer proofs than writing the full
  term by hand.
  \begin{lstlisting}[language=lean]
  import Mathlib

  theorem rat_le_trans : Transitive (· ≤ · : Rat → Rat → Prop) := by
    intro x y z hxy hyz
    exact Rat.le_trans hxy hyz
  \end{lstlisting}
  This proof performs the same steps, but is easier to read.
  Using \lstinline[language=lean]|by| we enter Lean's tactic mode.
  Move your cursor just after \lstinline[language=lean]|by|.
  The goal is initially displayed as
  \lstinline[language=lean]|⊢ Transitive fun x1 x2 ↦ x1 ≤ x2|.
  The tactic \lstinline[language=lean]|intro| introduces variables
  and hypotheses into the context (deconstructing universal
  quantifiers and implications).
  Now position your cursor just before
  \lstinline[language=lean]|exact| and observe the infoview again.
  The goal is now \lstinline[language=lean]|⊢ x ≤ z|, with the
  context showing the variables and hypotheses introduced by the
  previous tactic.
  The \lstinline[language=lean]|exact| tactic closes the goal by
  supplying a term that exactly matches it.
  You can hover over each tactic to see its definition and
  documentation.
\end{example}

% ADDED: transition
Having seen how Lean uses polymorphism and implicit arguments, let us
now look inside Mathlib to see how these lemmas are actually proved.

\section{Exploring Mathlib (The \lstinline[language=lean]|Rat|
  Structure)}\label{sec:exploring_mathlib}

In the previous examples we used predefined lemmas such as
\lstinline[language=lean]|Nat.le_trans| and
\lstinline[language=lean]|Rat.le_trans| to simplify the
presentation.
We can now dig into the implementation of these lemmas.
Let us look at the source code of
\lstinline[language=lean]|Rat.le_trans|.
The Mathlib~4 documentation website is at
\url{https://leanprover-community.github.io/mathlib4_docs}, and the
documentation for \lstinline[language=lean]|Rat.le_trans| is at
\url{https://leanprover-community.github.io/mathlib4_docs/Mathlib/Algebra/Order/Ring/Unbundled/Rat.html#Rat.le_trans}.
Click the ``source'' link there to jump to the implementation in the
Mathlib repository.
In editors like VS~Code, you can also jump directly to the definition
(Ctrl+click; Cmd+click on macOS).
Another way to inspect source code is by using
\lstinline[language=lean]|#print Rat.le_trans|.

\begin{lstlisting}[language=lean]
variable (a b c : Rat)
lemma le_trans (hab : a ≤ b) (hbc : b ≤ c) : a ≤ c := by
  rw [Rat.le_iff_sub_nonneg] at hab hbc
  have := Rat.add_nonneg hab hbc
  simp_rw [sub_eq_add_neg, add_left_comm (b + -a) c (-b),
    add_comm (b + -a) (-b), add_left_comm (-b) b (-a),
    add_comm (-b) (-a), add_neg_cancel_comm_assoc,
    ← sub_eq_add_neg] at this
  rwa [Rat.le_iff_sub_nonneg]
\end{lstlisting}

% REVISED: P27 — fixed misconception about rw and definitional equality.
The proof uses several tactics and lemmas from Mathlib.
The \lstinline[language=lean]|rw| (or \lstinline[language=lean]|rewrite|)
tactic is very common and mirrors the mathematical practice of
rewriting using equalities.
Given a proof \lstinline[language=lean]|h : A = B| (or an
iff \lstinline[language=lean]|h : A ↔ B|), the tactic
\lstinline[language=lean]|rw [h]| replaces occurrences of
\lstinline[language=lean]|A| with \lstinline[language=lean]|B| in the
goal.
Using \lstinline[language=lean]|rw [...] at hyp| performs the
rewrite in a specified hypothesis instead.
In this case, we rewrite the hypotheses
\lstinline[language=lean]|hab| and \lstinline[language=lean]|hbc|
via the Mathlib lemma
\lstinline[language=lean]|Rat.le_iff_sub_nonneg|, which states that
for any two rational numbers \lstinline[language=lean]|x| and
\lstinline[language=lean]|y|,
\lstinline[language=lean]|x ≤ y| is equivalent to
\lstinline[language=lean]|0 ≤ y - x|.
After rewriting, we have:
\begin{lstlisting}[language=lean]
  hab : 0 ≤ b - a
  hbc : 0 ≤ c - b
\end{lstlisting}

The \lstinline[language=lean]|have| tactic introduces an intermediate
result.
If you omit a name, Lean assigns it the default name
\lstinline[language=lean]|this|.
From \lstinline[language=lean]|hab| and \lstinline[language=lean]|hbc|
we derive that \lstinline[language=lean]|b - a| and
\lstinline[language=lean]|c - b| are nonnegative, hence their sum is
nonnegative:
\begin{lstlisting}[language=lean]
  this : 0 ≤ b - a + (c - b)
\end{lstlisting}

% REVISED: P27 — removed second incorrect claim about rw and def. equality.
The most involved step uses \lstinline[language=lean]|simp_rw| to
simplify the expression via a sequence of existing Mathlib lemmas.
The tactic \lstinline[language=lean]|simp_rw| is a variant of
\lstinline[language=lean]|rw| that can rewrite inside binders
(such as universal quantifiers $\forall x,\, P(x)$), whereas plain
\lstinline[language=lean]|rw| cannot.
A related tactic is \lstinline[language=lean]|simp|, which is more
powerful: it automatically searches through a database of lemmas
tagged as \textbf{simp lemmas} and applies them repeatedly to
simplify the goal.
After these simplifications we obtain:
\begin{lstlisting}[language=lean]
  this : 0 ≤ c - a
\end{lstlisting}

Clearly, the proof relies mostly on
\lstinline[language=lean]|Rat.add_nonneg|.
Mathlib defines \lstinline[language=lean]|Rat| as an instance of a
linear ordered field, implemented via a normalized fraction
representation: a pair of integers (numerator and denominator) with
positive denominator and coprime components~(\cite{mathlibdoc}).
To achieve this, it uses a \textbf{structure}.
In Lean, a structure is a dependent record type used to group
together related fields as a single data type.
Unlike ordinary records, the type of later fields may depend on the
values of earlier ones.
Defining a structure automatically introduces a constructor (usually
\lstinline[language=lean]|mk|) and projection functions that retrieve
the values of its fields.
Structures may also include proofs expressing properties that the
fields must satisfy.

\begin{lstlisting}[language=lean]
structure Rat where
  mk' ::
  num : Int
  den : Nat := 1
  den_nz : den ≠ 0 := by decide
  reduced : num.natAbs.Coprime den := by decide
\end{lstlisting}

To create a rational number in Mathlib, we use the
\lstinline[language=lean]|Rat.mk'| constructor, providing a
numerator and denominator.
The fields \lstinline[language=lean]|den_nz| and
\lstinline[language=lean]|reduced| are proofs that the denominator is
nonzero and that the numerator and denominator are coprime,
respectively.
% FIXED: P28 — broken sentence "decide tactic, which can We will"
These proofs are automatically generated by Lean's
\lstinline[language=lean]|decide| tactic, which can evaluate
decidable propositions by computation.
We will not delve into the type-theoretic formalization of
decidability, as it lies somewhat outside the scope of standard
mathematical practice.
Roughly speaking, a proposition is decidable if Lean can run a
computation that reduces the statement to
\lstinline[language=lean]|true| or \lstinline[language=lean]|false|.
We refer the interested reader to
\cite[Section~21.4]{nordstrom1990programming} for a deeper treatment.

\begin{example}
  Here is how we can define rational numbers in Lean:
  \begin{lstlisting}[language=lean]
    def half : Rat := Rat.mk' 1 2  -- or (1 / 2 : Rat)
    def third : Rat := Rat.mk' 1 3 -- or (1 / 3 : Rat)
  \end{lstlisting}
\end{example}

When working with rational numbers, or more generally with
structures, we must provide the required proofs as arguments to the
constructor (or Lean must be able to verify them automatically).
For instance, \lstinline[language=lean]|Rat.mk' 1 0| or
\lstinline[language=lean]|Rat.mk' 2 6| would be rejected.
In the case of rationals, Mathlib unfolds the definition through
\lstinline[language=lean]|Rat.numDenCasesOn|.
This principle states that, to prove a property of an arbitrary
rational number, it suffices to consider numbers of the form
$n / d$ with $d > 0$ and $\gcd(n, d) = 1$.
This reduction allows Mathlib to transform proofs about
\lstinline[language=lean]|ℚ| into proofs about
\lstinline[language=lean]|ℤ| and \lstinline[language=lean]|ℕ|, and
then lift the result back to rationals.

Let us return to \lstinline[language=lean]|Rat.add_nonneg|, which was
the important lemma used in the proof of
\lstinline[language=lean]|Rat.le_trans|.
We are going to construct our own simplified implementation of
rational numbers, inspired by Mathlib's approach, and use it to
prove the theorem.
We follow the same convention used in Mathlib: project to natural
numbers and integers first, then lift the result back to rationals.

Let us start by creating the structure:
\begin{lstlisting}[language=lean]
import Mathlib

structure myPreRat where
  num : Int
  den : Nat
  den_pos : 0 < den
\end{lstlisting}
Notice the similarity with Mathlib's definition.
We do not include the coprimality condition---the name
\lstinline[language=lean]|myPreRat| (for ``pre-rational'') will
become clear in Section~\ref{sec:quotients}, where we construct
proper rational numbers using quotient types.

% REVISED: P35 — reordered to present definitions before code that uses them.
Our goal is to prove
\lstinline[language=lean]|myPreRat.add_nonneg|.
Before we can state it, we need to define the operations
\lstinline[language=lean]|≤|, \lstinline[language=lean]|+|, and
\lstinline[language=lean]|0| for \lstinline[language=lean]|myPreRat|.
Operations such as addition or comparison need to be defined for each
type.
This is achieved through \textbf{type classes}.
Type classes provide a powerful and flexible way to specify properties
and operations that can be shared across different types, a mechanism
called \textbf{ad hoc
  polymorphism}~(\cite{wadler_blott_ad_hoc_polymorphism_1988}).
A standard example is overloaded multiplication: the same symbol
\lstinline[language=lean]|*| denotes multiplication of integers
(e.g., \lstinline[language=lean]|3 * 3|) and of floating-point
numbers (e.g., \lstinline[language=lean]|3.14 * 3.14|).
Lean exposes type classes for common operations:

\begin{lstlisting}[language=lean]
class Add (α : Type u) where
  add : α → α → α

class LE (α : Type u) where
  le : α → α → Prop
\end{lstlisting}

Type classes are, under the hood, structures whose fields describe
operations.
The important additional features of type classes are type inference
and \textbf{instances}.
\newpage
When we register a type as an instance of a type class, Lean will
automatically apply the operations defined by the class.
For example, we can register addition for naturals or integers:

\begin{lstlisting}[language=lean]
instance : Add Nat where
  add := Nat.add

instance : Add Int where
  add := Int.add
\end{lstlisting}

We now define instances for \lstinline[language=lean]|myPreRat|:
\begin{lstlisting}[language=lean]
instance : LE myPreRat where
  le r₁ r₂ := r₁.num * ↑r₂.den ≤ r₂.num * ↑r₁.den

instance : Add myPreRat where
  add r₁ r₂ := {
    num := r₁.num * ↑r₂.den + r₂.num * ↑r₁.den,
    den := r₁.den * r₂.den,
    den_pos := Nat.mul_pos r₁.den_pos r₂.den_pos
  }

def zero : myPreRat := { num := 0, den := 1, den_pos := by decide }

instance : OfNat myPreRat 0 where
  ofNat := zero
\end{lstlisting}

Once these instances are defined, Lean automatically infers which
operation to use when we write \lstinline[language=lean]|a + b| or
\lstinline[language=lean]|a ≤ b| for values of type
\lstinline[language=lean]|myPreRat|.
The \lstinline[language=lean]|↑| symbol denotes a coercion from
\lstinline[language=lean]|Nat| to \lstinline[language=lean]|Int|
(explained in Section~\ref{sec:coercions}).
With \lstinline[language=lean]|OfNat| we tell Lean that, in a context
expecting \lstinline[language=lean]|myPreRat|, the literal
\lstinline[language=lean]|0| should be interpreted as our
\lstinline[language=lean]|zero| definition.

We can now state and prove our lemmas.
We first need a helper lemma relating the ordering to the numerator:

% REVISED: P31 — simplified nonneg_iff as professor suggested
\begin{lstlisting}[language=lean]
namespace myPreRat

lemma nonneg_iff (r : myPreRat) : 0 ≤ r ↔ 0 ≤ r.num := by
  change 0 * r.den ≤ r.num * 1 ↔ _
  simp
\end{lstlisting}

This states that a pre-rational number is nonnegative if and only if
its numerator is nonnegative.
The \lstinline[language=lean]|change| tactic unfolds our definition
of \lstinline[language=lean]|≤| for
\lstinline[language=lean]|myPreRat|: since
\lstinline[language=lean]|0| has
\lstinline[language=lean]|num = 0| and
\lstinline[language=lean]|den = 1|, the inequality
\lstinline[language=lean]|0 ≤ r| becomes
\lstinline[language=lean]|0 * r.den ≤ r.num * 1|.
Then \lstinline[language=lean]|simp| simplifies the arithmetic to
obtain the equivalence with
\lstinline[language=lean]|0 ≤ r.num|.

We can now prove \lstinline[language=lean]|add_nonneg|:
\begin{lstlisting}[language=lean]
lemma add_nonneg (a b : myPreRat) : 0 ≤ a → 0 ≤ b → 0 ≤ a + b := by
  simp only [nonneg_iff]
  intro ha hb
  apply Int.add_nonneg
  · exact Int.mul_nonneg ha (Int.natCast_nonneg b.den)
  · exact Int.mul_nonneg hb (Int.natCast_nonneg a.den)

end myPreRat
\end{lstlisting}

The \lstinline[language=lean]|namespace| keyword defines a
self-contained module.
Outside of its scope, one refers to
\lstinline[language=lean]|add_nonneg| as
\lstinline[language=lean]|myPreRat.add_nonneg|.

We first simplify using \lstinline[language=lean]|nonneg_iff|, which
transforms the goal to
\lstinline[language=lean]|⊢ 0 ≤ a.num → 0 ≤ b.num → 0 ≤ (a + b).num|.
We then introduce the two hypotheses
\lstinline[language=lean]|ha : 0 ≤ a.num| and
\lstinline[language=lean]|hb : 0 ≤ b.num|.
By our definition of addition, the numerator of
\lstinline[language=lean]|a + b| is
\lstinline[language=lean]|a.num * ↑b.den + b.num * ↑a.den|.
Since the numerator is an integer, we can use lemmas from Mathlib.
We \lstinline[language=lean]|apply| the lemma
\lstinline[language=lean]|Int.add_nonneg|, which states that the sum
of two nonnegative integers is nonnegative.
The \lstinline[language=lean]|apply| tactic works backwards: given a
goal \lstinline[language=lean]|⊢ G| and a lemma
\lstinline[language=lean]|lemma : P → Q → G|, it replaces the goal
with two new subgoals \lstinline[language=lean]|⊢ P| and
\lstinline[language=lean]|⊢ Q|.
We close the first subgoal with
\lstinline[language=lean]|Int.mul_nonneg ha (Int.natCast_nonneg b.den)|,
where \lstinline[language=lean]|ha| provides the nonnegativity of
\lstinline[language=lean]|a.num|, and
\lstinline[language=lean]|Int.natCast_nonneg b.den| provides the
nonnegativity of the coercion of \lstinline[language=lean]|b.den|
from \lstinline[language=lean]|Nat| to \lstinline[language=lean]|Int|.
The second subgoal follows symmetrically.

Here is the full proof:
[\href{https://live.lean-lang.org/#codez=JYWwDg9gTgLgBAWQIYwBYBtgCMBQODOMUArgMYzFQCmcIAngArUBKKcA7qldTnHAHbEQcAFxwAkvxi84AEyr9RcAHIoZ8/gH1I+JQAY4AHjkK8/JCCr4wSUjXpMqraTmD9CSfnaUAZAKK0jCxsnNxUMug0UICBBHBQgEEEogC8cdEAdILCAFRwgImECWkacIAmRHHxGUJwOfnpGnhuHl40YgCCsrKBjs4cXDx8SO2pZclwAN4yfJkjMRXZeQVFANRls1XztQoANBMmiiIpM0U5C1s7GtoQuvsqKGkgxOgXuocKTyvnOjIAvnjyAGZwABe3AgSgcwXg11GAkq1z0m12IwAjAiPpcRlg6CZSMB5HAfg0YJ5vGIAPJ/VTwcFONgGUJ9OAQClsa7AqAQPCREAgJACCD8fhUADmmmAfwBAAooGCgjSYABKfQlOJwQAphHADKUoKtrpiZKR+YQSORoEYANwAPjgbiIoNQMgA7XBSKhPEKaAZjoUFMrtVMcki4GxUGa4PhQGAg/AQ3AqAAPWzRx3O138d0atbaopa1YB0Ph8Ch+OJuD2nBcnlB9qafj8wVCuAS3lYGVdFCKsSaoNwQBJhBnSi2+13ecsW7q6DIC5H+egsQBtWsC4Wi8UAXRkNvZpd5qFw/TAYFnEikaQGshrdeFyeL5GPMDuDwvS4brsbknv5hgAGEkIQn/W4Cwb1+Hla8E1vd8H0eRcAN3N8T0/H8/xg4Ug2A0CgA}{link to Lean live}]

% ADDED: transition to coercions
In the proof above, we repeatedly used the \lstinline[language=lean]|↑|
symbol to convert between \lstinline[language=lean]|Nat| and
\lstinline[language=lean]|Int|.
Let us now explain this mechanism properly.

\section{Coercions and Type Casting}\label{sec:coercions}

% REVISED: P32 — clarified the muddled discussion
When working with multiple numeric types such as
\lstinline[language=lean]|ℕ|, \lstinline[language=lean]|ℤ|, and
\lstinline[language=lean]|ℚ|, we often need to translate values
between them~(\cite{lewis_madelaine_simplifying_casts_coercions_2020}).
Lean provides two related mechanisms:

\begin{itemize}
  \item \textbf{Casting} refers to the explicit conversion of a
        value from one type to another, typically using functions like
        \lstinline[language=lean]|Nat.cast| or
        \lstinline[language=lean]|Int.cast|.
        These functions come with accompanying lemmas that preserve
        properties across type conversions, such as
        \lstinline[language=lean]|Nat.cast_pos| (if $n > 0$ as a
        natural number, then its image is positive as an integer).
  \item \textbf{Coercion} is a mechanism by which Lean
        \emph{automatically} inserts a cast when a type mismatch would
        otherwise occur.
        For example, if \lstinline[language=lean]|x : ℕ| and
        \lstinline[language=lean]|y : ℤ|, then in the expression
        \lstinline[language=lean]|x + y|, Lean automatically coerces
        \lstinline[language=lean]|x| to \lstinline[language=lean]|ℤ|.
\end{itemize}

The notation \lstinline[language=lean]|↑| makes a coercion explicit
in the source code, serving as a visual marker that a type conversion
is taking place.
For instance, in our \lstinline[language=lean]|LE| instance for
\lstinline[language=lean]|myPreRat|, the expression
\lstinline[language=lean]|↑r₂.den| coerces the denominator from
\lstinline[language=lean]|Nat| to \lstinline[language=lean]|Int| so
that it can be multiplied with the numerator.
\newpage
% REVISED: P32 — added concrete norm_cast and push_cast examples
Expressions involving coercions can become cumbersome.
To illustrate, consider
\lstinline[language=lean]|↑m + ↑n < (10 : ℤ)|, where
\lstinline[language=lean]|m n : ℕ| are cast to
\lstinline[language=lean]|ℤ|.
The expected simplified form is
\lstinline[language=lean]|m + n < 10| as an inequality over
\lstinline[language=lean]|ℕ|, since \lstinline[language=lean]|+|,
\lstinline[language=lean]{<}, and the numeral
\lstinline[language=lean]|10| are all polymorphic.

Lean provides the tactic family \lstinline[language=lean]|norm_cast|
to perform these simplifications automatically.
The tactic \lstinline[language=lean]|norm_cast| normalizes casts by
pushing them outward and eliminating redundant coercions:

\begin{lstlisting}[language=lean]
example (m n : ℕ) (h : ↑m + ↑n < (10 : ℤ)) : m + n < 10 := by
  exact_mod_cast h
\end{lstlisting}

Here \lstinline[language=lean]|exact_mod_cast| is a variant of
\lstinline[language=lean]|exact| that automatically handles coercion
simplification.
A related tactic, \lstinline[language=lean]|push_cast|, pushes casts
\emph{inward} toward the leaves of an expression, which is useful
when one wants to reason within a single type:

\begin{lstlisting}[language=lean]
example (m n : ℕ) : (↑(m + n) : ℤ) = ↑m + ↑n := by
  push_cast
\end{lstlisting}

% UPDATED: chapter-closing transition
Having established how Lean handles type conversions, we now have
all the tools needed to address a fundamental problem with
\lstinline[language=lean]|myPreRat|: different representations of
the same rational number are not identified as equal.
In the next chapter, we resolve this using quotient types and
discuss how Mathlib organizes its algebraic hierarchy using
type classes.